\section{Safety, Accessibility, and External Validity}

AR blocks are restricted to a caregiver-prepared indoor area with sufficient free space, stable lighting, and no immediate fall or collision hazards. The application does not require running, outdoor geolocation, road-adjacent play, stairs, or movement while visual attention is captured by the display.

The child must have immediate pause and stop controls. Dizziness, nausea, visual discomfort, headache, repeated tracking loss, unsafe movement, or an explicit request to stop terminates the current AR block according to a predefined safety pathway.

Movement tasks should offer seated or reduced-range alternatives where possible. Children with common co-occurring motor or sensory difficulties should not be excluded merely because the original interaction can be adapted safely. The adapted condition and its consequences for interpretation must be recorded.

Because the programme uses private phones, device compatibility and hardware-related exclusion are measured directly rather than removed through standardized handset provision. The AR module therefore reports qualification failure, fallback or exclusion by coarse device-capability class and distinguishes technical tracking/runtime failure from participant performance.

Home environments differ in available floor area, visual texture, lighting, furniture, caregiver support, and network conditions. These variables should be recorded because they limit external validity and may moderate feasibility. They must not be treated as random noise that disappears through a larger sample.

Rear-camera environmental imagery is processed locally and omitted from standard telemetry. Spatial maps or anchors that could reveal the home environment should be retained only when scientifically necessary and covered by an explicit data-governance plan.
